\scanpage{55}
\mspage{48}

\deleted{\unsure{on}} $s\backslash I'$ est discrète, on a donc l'isom. fonctoriel en $X'$ dans $M(I')$

\[ \alpha_s^*(u_!(X'))\simeq\coprod_{\substack{(s'\unsure{,|;}\unsure{h|\lambda}:s\to u(s'))\\\text{dans }\mathrm{Ob}\,s\backslash I'}}X(s'), \]

\deleted{\unsure{donc}} donc \inserted{si $f:X'\unsure{-|\to}Y'$,} $\alpha_s^*u_!(f)$ s'explicite comme une somme directe \inserted{dans $\mathrm{Fl}(M)$} de flèches de la forme

\[ f_{s'}:X'(s')\longrightarrow Y'(s'), \]

lesquelles sont dans $C$ resp. $\mathrm{TC}$ puisque $f$ est dans $C(I')$ resp. $\mathrm{TC}(I')$. Comme $C,\mathrm{TC}$ sont stables par sommes directes quelconques. ($C,\mathrm{TC}$ étant clos à gauche), on gagne.

\uline{Corollaire 2} Soit $I$ dans $\mathrm{Cat}$, \deleted{Alors \unsure{pour $\forall\,s$}} $f\in\mathrm{Fl}(M(I))$. Pour que $f$ soit dans $F(I)$ (resp. \deleted{dans} $\mathrm{TF}(I)$), il \uline{faut} que $\forall s\in\mathrm{Ob}\,I$, $f_s:X(s)-Y(s)$ le soit. (Mais cette condition n'a pas \unsure{du|un} \deleted{\unsure{tout}} l'air d'être suffisante !) En particulier, soit $X\in\mathrm{Ob}\,M(I)$, alors $X\in\mathrm{Ob}\,M(I)_f$ $\Longrightarrow X(s)\in\mathrm{ob}\,M_f\quad\forall s\in\mathrm{Ob}\,I$.

On applique le cor. 1 à $\alpha_{\unsure{e|s}}:e-I$.

\uline{Scholie} La prop. 1 \inserted{(\unsure{relations|relation} (1))} montre que le foncteur $u^*$ transforme carrés \inserted{$W$-}cocartésiens en \unsure{itou|idem}, suite exacte longue de suspensions en \unsure{itou|idem} -- c'est ``l'exactitude à droite ``homotopique'' de $u^*$. On s'attend à ce qu'il en soit de $=$ pour les carrés cartésiens et les suites exactes
% archive sheet number: 54
